% !TEX root = ../algebra_02.tex

\section{Факторгруппа}

\begin{Definition}{Фактормножество}{def:quotient-set}
	Пусть на множестве $X$ задано отношение эквивалентности $\sim$. Тогда \emph{фактормножество} $X/{\sim}$ определяется как множество всех классов эквивалентности:
	\[
	X/{\sim}:=\{[x] \mid x\in X\}.
	\]
\end{Definition}

\begin{Remark}{Фактормножество группы по нормальной подгруппе}{rem:group-quotient-set}
	Пусть $G$ --- группа, $H \triangleleft G$. Тогда на $G$ задано отношение эквивалентности
	\[
	g_1\sim g_2 \Longleftrightarrow g_1^{-1}g_2\in H.
	\]
	Его классы эквивалентности --- это левые смежные классы по $H$, поэтому соответствующее фактормножество имеет вид
	\[
	G/H:=\{gH\mid g\in G\}.
	\]
	На $G/H$ задаётся операция
	\[
	(gH)\cdot(g'H):=(gg')H.
	\]
	Так как $H$ нормальна, левые и правые смежные классы совпадают: $gH=Hg$ для всех $g\in G$.
\end{Remark}

\begin{Definition}{Факторгруппа}{def:quotient-group}
	Пусть $G$ — группа, $H \triangleleft G$.

	Множество левых смежных классов $G/H$ с операцией $(gH) \cdot (g'H) = (gg')H$ называется \emph{факторгруппой} группы $G$ по нормальной подгруппе $H$.
\end{Definition}

\begin{Proposition}{Корректность умножения смежных классов}{prop:well-defined}
	Если $H \triangleleft G$, то операция
	\[
	(gH)(g'H) := (gg')H
	\]
	корректно определена, то есть не зависит от выбора представителей смежных классов.
\end{Proposition}

\begin{proof}
	Пусть $g_1H=g_2H, \qquad g_1'H=g_2'H$.

	Тогда существуют $h,h'\in H$ такие, что $g_2=g_1h, \qquad g_2'=g_1'h'$. Следовательно, $g_2g_2' = g_1 h g_1' h' = g_1 g_1'\bigl((g_1')^{-1} h g_1'\bigr)h'$. Так как $H \triangleleft G$, имеем $(g_1')^{-1} h g_1' \in H$. Поэтому $\bigl((g_1')^{-1} h g_1'\bigr)h' \in H,$ и значит, $g_2g_2'H = g_1g_1'H.$

	Следовательно, $(g_2H)(g_2'H)=(g_1H)(g_1'H).$

	Корректность доказана.
\end{proof}

\begin{Theorem}{Групповая структура на $G/H$}{thm:quotient-group}
	Если $H \triangleleft G$, то множество $G/H$ с операцией
	\[
	(gH)(g'H)=(gg')H
	\]
	является группой.
\end{Theorem}

\begin{proof}
	Проверим аксиомы группы.
	\begin{enumerate}
		\item Ассоциативность:
		\[
		\bigl((g_1H)(g_2H)\bigr)(g_3H) = (g_1g_2g_3)H = (g_1H)\bigl((g_2H)(g_3H)\bigr).
		\]
		\item Нейтральный элемент: $eH=H$, и для любого $g\in G$
		\[
		(eH)(gH)=gH=(gH)(eH).
		\]
		\item Обратимый элемент: для $gH\in G/H$
		\[
		(gH)(g^{-1}H)=eH=(g^{-1}H)(gH).
		\]
	\end{enumerate}
	Следовательно, $G/H$ --- группа.
\end{proof}

По определению эта группа называется \emph{факторгруппой} группы $G$ по нормальной подгруппе $H$.

\begin{Remark}{}{rem:not-subgroup}
	Вообще говоря, $G/H$ не является подмножеством группы $G$. Это новая группа, элементы которой --- смежные классы. Исключения возникают лишь в тривиальных случаях $H=\{e\}$ и $H=G$.
\end{Remark}

\begin{Example}{Простейшие примеры факторгрупп}{ex:trivial-quotients}
	Имеют место изоморфизмы
	\[
	G/\{e\} \cong G, \qquad G/G \cong \{e\}.
	\]
	Если группа записана аддитивно, то факторгруппа тоже записывается аддитивно. Например,
	\[
	\ZZ/m\ZZ = \{[a]_m \mid a\in \ZZ\}, \qquad [a]_m + [b]_m := [a+b]_m.
	\]
\end{Example}
